\documentclass{article}
\usepackage{amsmath,amssymb}
\usepackage{booktabs}
\usepackage{hyperref}
\usepackage{geometry}
\geometry{margin=1in}

\newcommand{\pos}[1]{\left(#1\right)^{+}}
\setlength{\emergencystretch}{3em}
\title{The Cost of Wordplay: an Educational Programming Platform}
\author{Issy Chen, Siran Luo}
\date{}

\begin{document}
\maketitle

\begin{abstract}
Wordplay is a free coding platform built by Professor Amy Ko at UW for students, teachers, and families. Running it costs money from Firebase, Claude AI tools, and fixed items like backups. We build a daily cost model. Student activity creates database operations, and each day we subtract the free usage limit before counting charges. Monthly cost is the total of the daily costs plus the fixed costs. We build the model in steps, starting with the biggest cost parts and then adding the rest. We use Monte Carlo simulation to try different class sizes and activity levels, and we check which factors change the cost the most.
\end{abstract}

\section{Problem Description}

\paragraph{Background.}
Wordplay is an open-source coding platform built with educational justice in mind. It serves students and children, and also families, teachers, and the general public. Keeping it online costs real money. User projects are saved and synced through Firebase (Hosting, Firestore, Authentication, and Functions). Firebase charges for Firestore reads, writes, and deletes, and also for storage, function calls, outgoing data, and hosting. Wordplay also uses Claude AI to help students translate projects and chats, and to check how easy text is to read. Some costs are fixed, which means we pay them even if nobody uses the app. Examples are scheduled jobs, backups, and deploys. Firebase also gives free usage limits that reset every day (or month). So the final bill depends not only on how much people use Wordplay, but also on whether they go over these free limits.

\paragraph{Goal.}
The goal of this model is to predict the cost of Wordplay for some number of users to identify drivers of costs, given the fixed costs per month, the bounds of the daily cap per user, and the amount of function costs.

\paragraph{Why it matters.}
Wordplay is run by a research group at UW and is meant to stay free and easy for students to use. Some costs do not grow in a simple way with the number of users. For example, the number of reads in a shared gallery is
\[ R_{\text{gallery}} = N\,S\,(N-1)\,H, \]
where $N$ is the number of students, $S$ is how many times a student saves per hour, and $H$ is the hours of editing. This grows with $N^2$, so a few large classes can make the cost jump much higher than a simple straight-line guess. Other costs do not depend on users at all. For example, keeping backups of the database costs about five times as much as storing the data. A good cost forecast helps the team plan their budget, set reasonable daily limits for each user, and decide which technical changes would save the most money. It also helps the server side run smoothly, so the platform is easier to maintain.

\paragraph{Questions we want to answer.}
\begin{enumerate}
    \item What happens to the monthly bill when we go from $N_1$ students up to $N_2$ students, and when do we start going over the free usage limits?
    \item Which part costs the most: Firestore reads, writes, or deletes, Claude AI requests, hosting, storage, or the fixed costs?
    \item How does the size of a class or shared gallery change the total cost, since gallery reads grow with $N^2$?
    \item Within the bounds of the daily cap per user, which cap keeps the monthly cost below our budget $B$?
    \item How much does the cost change if our guesses for numbers like writes per student per hour $w_h$ and saves per student per hour $S$ are off?
\end{enumerate}

\section{Simplifications}

The real Wordplay platform is complicated, and its cost depends on many things we cannot track exactly. To get a model we can solve and understand, we make the following simplifications.

\begin{enumerate}
    \item \textbf{We build the model in steps.}
    Our first step includes the biggest cost parts: the Firestore reads, writes, and deletes from student activity, the Claude AI requests, and the fixed monthly costs. In later steps we add Firestore storage, function calls, outgoing data, hosting, App Check (reCAPTCHA), and a few extra items. The full model follows the total in the cost sheet:
    \[
    C_{\text{total}} = C_{\text{Claude}} + C_{\text{Firebase}} + C_{\text{fixed}} + C_{\text{App Check}} + C_{\text{getWebpage}} + C_{\text{dev}} + C_{\text{one-off}},
    \]
    where
    \[
    C_{\text{Firebase}} = C_{\text{Firestore}} + C_{\text{storage}} + C_{\text{calls}} + C_{\text{outbound}} + C_{\text{hosting}} + C_{\text{hosting storage}}.
    \]
    Building in steps lets us understand every piece we add. The sensitivity analysis will also show which pieces matter most.

    \item \textbf{We model one month, one day at a time.}
    Firebase's free usage limits reset every day and do not carry over, so the cost of one day does not depend on the days before it. This lets us compute the cost of each day and add them up for the month:
    \[
    C_{\text{Firestore}} = \sum_{t=1}^{\text{days}} C_{\text{day}}(t), \qquad
    C_{\text{day}} = \frac{0.09\,(W-20000)^+ + 0.03\,(R-50000)^+ + 0.01\,(D-20000)^+}{10^5},
    \]
    where $W$, $R$, $D$ are the writes, reads, and deletes in one day, and $(x)^+ = \max(x,0)$, so usage inside the free limit costs nothing. We treat all days as equally busy, so we miss busy or slow times like assignment deadlines or summer break.

    \item \textbf{Our activity numbers are estimates.}
    We do not have measured data on how students use the platform, so we use the estimates from the cost sheet: $w_h \approx 2500$ writes per student per hour (720 presence $+$ 1500 typing $+$ 300 saves), deletes $D \approx 1500\,NH$, and Claude token counts $U=1500$, $W_c=R_c=2000$, $O+T=5000$. The model is only as good as these inputs, so if they are off, the predicted cost is off too. To deal with this, we run a sensitivity analysis where we change each estimate and see how much the total cost moves. Real user logs would let us use real numbers instead.

    \item \textbf{We assume every student in a class acts the same way.}
    All students have the same editing hours $H$ and saves per hour $S$. Students in one class usually work on the same assignment at the same time, so this is a reasonable first step. It also lets us describe a class with just a few numbers: $N$, $H$, and $S$. In real life, some students and classes are much more active than others, so our estimate could be too low for busy classes.

    \item \textbf{We use simple formulas for shared work and Claude requests.}
    Reads from shared work are hard to count exactly, so we use the approximations from the cost sheet: $R_{\text{gallery}} = NS(N-1)H$ for a gallery of $N$ students, and $R_{\text{project}} \approx w_p\,(V-1)$ for a project edited together, with at most $E \le 4$ editors. For Claude, checking the reading level of $n$ text strings is split into $k=\lceil n/40\rceil$ chunks of up to 40 strings. Since the chunks are about the same size, we use one average cost, $C_{\text{check}} \approx 0.116 + 0.106\,(k-1)$ dollars. The real cost changes with how long each prompt and reply is.

    \item \textbf{Fixed costs are constants, and daily caps are inputs.}
    Our goal is to predict cost given the fixed costs per month and the bounds of the daily cap per user. The fixed costs are the items we pay whether or not anyone uses the app: scheduled jobs, secrets, backups and point-in-time recovery, the container registry, deploys, and outside items like the email service and the domain name:
    \[
    C_{\text{fixed}} = C_{\text{scheduler}} + C_{\text{secrets}} + C_{\text{backups}} + C_{\text{PITR}} + C_{\text{registry}} + \sum_{\text{days}} C_{\text{build}} + C_{\text{Resend}} + C_{\text{domain}}.
    \]
    For example, with $j=8$ scheduled jobs per project and $P=2$ projects, $C_{\text{scheduler}} = 0.10\,(J-3)^+ = \$1.30$ per month. We assume each backup is the size of the database, as in the cost sheet, so backups and recovery cost about $0.93\,G$ per month for a database of $G$ GiB. For a database of a few GiB, the fixed costs add up to a few dollars a month, plus the email service and the domain. We treat these as constants for a given month. We treat the daily cap per user as a strict limit $q$ between a lower bound $q_{\min}$ and an upper bound $q_{\max}$, so we can ask which cap keeps the cost under a budget. We do not model what students do once they hit the cap.

    \item \textbf{Costs that do not depend on students are set as inputs.}
    Some items in the total do not depend on the number of students in a simple way, so we set them as inputs instead of simulating them. The App Check cost depends on the number of active devices $d$, using about $A \approx 4.3\,d$ assessments per month and the cost sheet's assumption that the \$8 fee applies once $A$ is over 10{,}000. The \texttt{getWebpage} function can be called by anyone, so we use $C_{\text{getWebpage}} \approx 0.12\,e_w$ for $e_w$ GB of outgoing data. Developers' Claude use, $C_{\text{dev}} \approx 3\ell + C_{\text{changelog}} + C_{\text{scripts}}$ for $\ell$ new languages a month, follows how much the team works, and the cost sheet says it was not measured. One-off costs, like restoring a backup ($0.20$ per GiB restored) or running a database migration, happen only sometimes. We try a few values for these items and see whether they change the answer.

    \item \textbf{Some costs are left out.}
    We ignore extra charges for Cloud Function CPU and memory above the monthly free amount (200{,}000 CPU-seconds and 400{,}000 GB-seconds), and Authentication fees above 50{,}000 monthly active users. We also leave out the items the cost sheet does not give prices for, such as logging, Firestore outgoing data above the free amount, security rule reads, crawlers, the email inbox, and tax. The model will underestimate the bill if usage gets high enough to trigger these charges.

    \item \textbf{Other assumptions.}
    We use the prices from the cost sheet (Claude API and Google Cloud us-central1, in US dollars) and assume they stay the same. We also choose the number of users $N$ ourselves instead of predicting how users join or leave, so our model answers ``how much would $N$ users cost?''
\end{enumerate}

\section{Mathematical Model}

All costs are in US dollars. Prices are for Claude's first-party API and
Google Cloud us-central1. Per-hour activity counts marked ``estimate'' are
assumptions, not measurements.

\paragraph{Notation.} The positive part of $x$ is
\[
\pos{x} = \max(x, 0),
\]
so usage inside a free tier costs nothing.

\subsection{Claude}

\subsubsection{Variables}

\begin{center}
\begin{tabular}{ll}
\toprule
Symbol & Meaning \\
\midrule
$I$ & Input tokens in a Haiku call \\
$O$ & Output tokens in a call \\
$U$ & Uncached input tokens in an Opus call \\
$W_c$ & Tokens written to the prompt cache \\
$R_c$ & Tokens read from the prompt cache \\
$T$ & Thinking tokens, billed as output \\
$n$ & Strings being checked \\
$k$ & Chunks of up to 40 strings \\
\bottomrule
\end{tabular}
\end{center}

\subsubsection{Translation: \texttt{getLLMTranslations} (Haiku 5.5)}

\[
C_{\text{translate}} = \frac{0.10\,I + 0.50\,O}{10^{6}}
\]

\subsubsection{Reading-level check: \texttt{analyzeLocalization} (Opus 5.5)}

\[
C_{\text{chunk}} =
\frac{4\,U + 5\,W_c + 0.20\,R_c + 20\,(O + T)}{10^{6}}
\]

\[
k = \left\lceil \frac{n}{40} \right\rceil
\]

With the estimates $U = 1500$, $W_c = R_c = 2000$ and $O + T = 5000$:

\[
C_{\text{check}} \approx 0.116 + 0.106\,(k - 1)
\]

\subsection{Firebase}

\subsubsection{Variables}

\begin{center}
\begin{tabular}{ll}
\toprule
Symbol & Meaning \\
\midrule
$W$ & Firestore document writes in a day** \\
$R$ & Firestore document reads in a day** \\
$D$ & Firestore document deletes in a day** \\
$G$ & Firestore storage in GiB \\
$N$ & Students in one gallery \\
$H$ & Hours of editing** \\
$S$ & Project saves per student per hour** \\
$E$ & Editors in a shared project, $E \le 4$ \\
$V$ & People with a shared project open \\
$w_h$ & Writes per student per hour** \\
$w_p$ & Writes to one shared project in the period \\
$c$ & Function calls in a month** \\
$e$ & Function outbound data in GB per month \\
$b$ & Hosting downloads in GB per day \\
$h$ & Hosting storage in GB \\
\multicolumn{2}{l}{**Indicates user behaviour} \\
\bottomrule
\end{tabular}
\end{center}

\subsubsection{Firestore cost per day}

\[
C_{\text{day}} =
\frac{0.09\,\pos{W - 20000}
    + 0.03\,\pos{R - 50000}
    + 0.01\,\pos{D - 20000}}{10^{5}}
\]

Free allowances reset daily and do not roll over, so monthly cost is the sum
of the daily costs:

\[
C_{\text{Firestore}} = \sum_{d=1}^{\text{days}} C_{\text{day}}(d)
\]

\subsubsection{Activity estimates}

\[
w_h = \underbrace{720}_{\text{presence}}
    + \underbrace{1500}_{\text{typing, estimate}}
    + \underbrace{300}_{\text{saves, estimate}}
    \approx 2500
\]

\[
D \approx 1500\,N H \quad \text{(typing updates removed by cleanup)}
\]

Reads in one shared gallery:

\[
R_{\text{gallery}} = N\,S\,(N - 1)\,H
\]

Reads in one project edited together:

\[
R_{\text{project}} \approx w_p\,(V - 1)
\]

\subsubsection{Class cost before the free tier}

With $w_h = 2500$ and $S = 300$:

\[
C_{\text{class}} = N H \bigl(0.00225 + 0.00009\,(N - 1)\bigr)
\]

\subsubsection{Other Firebase costs}

\begin{align*}
C_{\text{storage}} &= 0.15\,\pos{G - 1} && \text{per month} \\
C_{\text{calls}} &= \frac{0.40\,\pos{c - 2000000}}{10^{6}} && \text{per month} \\
C_{\text{outbound}} &= 0.12\,\pos{e - 5} && \text{per month} \\
C_{\text{hosting}} &= 0.15\,\pos{b - 0.36} && \text{per day} \\
C_{\text{hosting storage}} &= 0.026\,\pos{h - 10} && \text{per month}
\end{align*}

Not included: function CPU and memory past 200{,}000 CPU-seconds and
400{,}000 GB-seconds a month, and Authentication past 50{,}000 monthly active
users. Those rates were not checked.

\subsection{Fixed Costs}
\subsubsection{Cap Costs}
Many aspects within Wordplay operate within the free tier, fixed costs above the free tier is as follows. 

\subsubsection{What counts as fixed}

These are paid whether or not anyone uses the app. Most scale with the number
of Firebase projects (\texttt{wordplay-prod} and \texttt{wordplay-dev}) or with
the total size of the database, not with activity. Claude has no fixed fee.

\subsubsection{Variables}

\begin{center}
\begin{tabular}{ll}
\toprule
Symbol & Meaning \\
\midrule
$P$ & Firebase projects with the functions deployed \\
$J$ & Scheduled jobs on the billing account \\
$j$ & Scheduled jobs per project ($j = 8$) \\
$s$ & Active secret versions on the billing account \\
$a$ & Secret accesses in a month \\
$G$ & Firestore database size in GiB \\
$B_d$ & Daily backups kept ($B_d = 7$ in prod) \\
$B_w$ & Weekly backups kept ($B_w = 14$ in prod) \\
$G_r$ & Function container images stored, in GiB \\
$m$ & Cloud Build minutes in a day \\
\bottomrule
\end{tabular}
\end{center}

\subsubsection{Scheduled jobs (Cloud Scheduler)}

Each job costs \$0.10 per 31 days. The first 3 jobs on the billing account are
free, and executions are not billed.

\[
J = j\,P, \qquad
C_{\text{scheduler}} = 0.10\,\pos{J - 3} \quad \text{per month}
\]

With $j = 8$ and $P = 2$: $C_{\text{scheduler}} = 0.10 \times 13 = \$1.30$.

Running the jobs costs about 1{,}000 function calls per project each month
(mostly the hourly \texttt{compactProjectUpdates} and the 4-hourly chat
digest), well inside the free 2 million.

\subsubsection{Secrets (Secret Manager)}

\[
C_{\text{secrets}} = 0.06\,\pos{s - 6}
                   + \frac{0.03\,\pos{a - 10000}}{10^{4}}
                   \quad \text{per month}
\]

Three secrets in two projects is $s = 6$, which is free. A secret is read when
a function instance starts, so $a$ grows with cold starts.

\subsubsection{Backups and point-in-time recovery (prod only)}

Assuming each backup is the size of the database:

\[
C_{\text{backups}} = 0.03\,(B_d + B_w)\,G = 0.63\,G
\qquad
C_{\text{PITR}} = 0.15\,G
\quad \text{per month}
\]

Protecting the data costs about five times what storing it does. Ignoring the
free 1\,GiB, the total for data is:

\[
C_{\text{data}} \approx 0.15\,G + 0.63\,G + 0.15\,G = 0.93\,G
\quad \text{per month}
\]

These have no free tier and grow with total data kept, not with activity.

\subsubsection{Deploys}

Every push to \texttt{main} redeploys hosting, rules and every function.

\begin{align*}
C_{\text{build}} &= 0.003\,\pos{m - 120} && \text{per day} \\
C_{\text{registry}} &= 0.10\,\pos{G_r - 0.5} && \text{per month}
\end{align*}

$G_r$ depends on whether old images are cleaned up, which is set in Artifact
Registry rather than in the repository. Hosting files are about 223\,MB per
release, inside the free 10\,GB unless many old releases are kept.

\subsubsection{Outside Google Cloud}

\begin{center}
\begin{tabular}{ll}
\toprule
Item & Cost \\
\midrule
Resend plan & Not checked; free tier or a monthly plan \\
\texttt{wordplay.dev} domain & Yearly, from the registrar \\
GitHub & Free \\
Anthropic & No fixed fee; pay per token only \\
\bottomrule
\end{tabular}
\end{center}

\subsubsection{Total fixed cost}

\[
C_{\text{fixed}} =
C_{\text{scheduler}} + C_{\text{secrets}} + C_{\text{backups}}
+ C_{\text{PITR}} + C_{\text{registry}} + \sum_{\text{days}} C_{\text{build}}
+ C_{\text{Resend}} + C_{\text{domain}}
\]

For a database of a few GiB this is a few dollars a month, plus Resend and the
domain.

\subsection{Other costs}

\subsubsection{Variables}

\begin{center}
\begin{tabular}{ll}
\toprule
Symbol & Meaning \\
\midrule
$A$ & reCAPTCHA assessments in a month, across the organization \\
$d$ & Active devices that call a Cloud Function in a month \\
$e_w$ & Outbound GB from \texttt{getWebpage} in a month \\
$L$ & New languages translated in a month \\
$G_{\text{restore}}$ & GiB restored from a backup \\
$r_m$ & Documents read by a migration \\
$w_m$ & Documents written by a migration \\
\bottomrule
\end{tabular}
\end{center}

\subsubsection{App Check (reCAPTCHA)}

A token lasts 7 days and is only requested when a Cloud Function is called,
so each active device makes about one assessment a week:

\[
A \approx 4.3\,d
\]

\[
C_{\text{App Check}} =
\begin{cases}
0 & A \le 10000 \\
8 & 10000 < A \le 100000 \\
8 + \dfrac{A - 100000}{1000} & A > 100000
\end{cases}
\quad \text{per month}
\]

The pricing page lists \$1 per 1{,}000 above 100{,}000 assessments without
saying whether the \$8 fee still applies; this formula assumes it does.

\subsubsection{\texttt{getWebpage}}

The function fetches any URL with no sign-in, App Check or size limit, so its
cost is set by whoever calls it. Its outbound data shares the function
allowance of 5\,GB, so once that is used each extra GB costs:

\[
C_{\text{getWebpage}} \approx 0.12\,e_w
\quad \text{per month, plus running time}
\]

\subsubsection{Development use of Claude}

These use the project's Anthropic key from developers' machines and GitHub
Actions, and follow how much the team works rather than how many users there
are:

\[
C_{\text{dev}} \approx 3\,L + C_{\text{changelog}} + C_{\text{scripts}}
\quad \text{per month}
\]

$C_{\text{changelog}}$ is 29 translations per new changelog entry, and
$C_{\text{scripts}}$ covers occasional Opus scripts such as
\texttt{condense-changelog}. Neither was measured.

\subsubsection{One-off costs}

\begin{align*}
C_{\text{restore}} &= 0.20\,G_{\text{restore}} \\
C_{\text{migration}} &= \frac{0.03\,r_m + 0.09\,w_m}{10^{5}}
\end{align*}

A migration that reads and rewrites 100{,}000 documents costs about \$0.12.
Testing on \texttt{test.wordplay.dev} also uses real Firestore, function and
email costs, at the rates in Section~2.

\subsubsection{Not modeled}

\begin{center}
\begin{tabular}{ll}
\toprule
Item & Note \\
\midrule
Cloud Logging & Free allowance per project; price above it not checked \\
Firestore outbound data & Free up to 10\,GiB a month; price above it not checked \\
Security rule reads & One extra read when a rule checks a gallery document \\
Crawlers & Run \texttt{getPagePreview} and \texttt{getSitemap} like visitors \\
GitHub Actions & Free while the repository is public \\
Email inbox for \texttt{hi@} & Depends on the provider \\
Tax & Depends on billing location \\
\bottomrule
\end{tabular}
\end{center}

\subsubsection{Total}

\[
C_{\text{total}} =
C_{\text{Claude}} + C_{\text{Firebase}} + C_{\text{fixed}}
+ C_{\text{App Check}} + C_{\text{getWebpage}} + C_{\text{dev}}
+ C_{\text{one-off}}
\]

% Paste this section into your main document (after the Mathematical Model section).
% Needs: \usepackage{amsmath,amssymb}   (pseudocode uses plain verbatim, no extra packages)

\section{Solution of the Mathematical Problem}

\subsection{Overview}

Our model has a cost function (the formulas from the Mathematical Model section) and random student behavior. We solve it in three steps:

\begin{enumerate}
    \item \textbf{Deterministic baseline.} Every student acts the same way every day, and we compute the cost directly from the formulas.
    \item \textbf{Monte Carlo simulation.} We randomly generate many possible months of student activity, compute the cost of each one, and look at the average and the spread.
    \item \textbf{Analysis.} We use the simulation to find the cost drivers (sensitivity analysis) and to choose a daily cap $q$ that keeps the cost under the budget $B$.
\end{enumerate}

We will write the code in Python, using NumPy for random numbers and Pandas for tables of results.

\subsection{Step 1: Deterministic baseline}

In the baseline, each class meets for the same number of hours $H$ every school day, and every student makes the same number of translation requests. Then the counts $W_t$, $R_t$, $D_t$ are the same on every school day, and the monthly cost is the daily cost added up over the days of the month. No random numbers are needed, so this is a fast way to check the formulas.

\paragraph{Example: when do we go over the free limits?}
Take one class of $N$ students that edits for $H=1$ hour, with $S=300$ and $w_h=2500$. The free limits are 20{,}000 writes, 50{,}000 reads, and 20{,}000 deletes per day. We find the smallest class size that goes over each limit:
\begin{itemize}
    \item Writes: $2500\,N > 20000$ means $N \ge 9$.
    \item Reads: $300\,N(N-1) > 50000$ means $N(N-1) > 166.7$, so $N \ge 14$ (since $13\cdot 12 = 156$ is too small and $14 \cdot 13 = 182$ is enough).
    \item Deletes: $1500\,N > 20000$ means $N \ge 14$.
\end{itemize}
So a single class starts paying for writes at 9 students, and for reads and deletes at 14 students. This kind of calculation answers our question about which activity levels push the platform past its free limits.

\subsection{Step 2: Monte Carlo simulation}

\paragraph{Why we need it.}
The cost function has kinks at the free limits, and it bends upward after each one (it is a convex function). Because of this, computing the cost from the \emph{average} activity gives a number that is too small. Days with high activity cost more than days with low activity save. Here is a small example. Take one class of 30 students. Suppose the class either does not meet ($H=0$) or meets for 2 hours ($H=2$), each with probability $1/2$, so the average is $H=1$ hour.
\begin{itemize}
    \item Using the average, $H=1$: the daily cost is $\$0.115$ (from our example in the Mathematical Model section).
    \item With $H=2$: $W=150{,}000$, $R=522{,}000$, $D=90{,}000$, so the cost is $\dfrac{0.09(130000)+0.03(472000)+0.01(70000)}{10^5} = \dfrac{26560}{10^5} \approx \$0.266$. With $H=0$ the cost is \$0. The expected daily cost is $\tfrac12(0.266) + \tfrac12(0) \approx \$0.133$.
\end{itemize}
The expected cost, \$0.133, is about 15\% higher than the cost of the average day, \$0.115. This is why we simulate the random days instead of only using averages.

\paragraph{What we sample.}
In each simulated month, we randomly generate the following. These distributions are our own simple choices, since we do not have real usage logs, and we vary them in the sensitivity analysis.
\begin{itemize}
    \item Whether class $i$ meets on day $t$: it meets with probability $p_{\text{meet}}$ (for example, $p_{\text{meet}} = 2/3$, which gives about 20 school days in a 30-day month).
    \item If it meets, the hours $H_{i,t}$: for example, 1 or 2 hours with equal probability.
    \item The number of translation requests $X_{u,t}$ each student wants to make: a Poisson random variable with mean $\lambda$ (a standard choice for counts of random events).
    \item The number of strings $n_t$ checked for reading level: a Poisson random variable with mean $\mu$.
\end{itemize}
The class sizes $N_i$, the number of classes $K$, the cap $q$, and all the prices and fixed costs are inputs that we choose for each scenario.

\paragraph{Pseudocode for one simulated month.}
\begin{verbatim}
function SimulateMonth(inputs):
    total_Firestore = 0, total_Claude = 0, total_calls = 0
    for t = 1 to days:
        W = 0, R = 0, D = 0
        for i = 1 to K:
            if Uniform(0,1) < p_meet:
                H = sample hours
                W = W + w_h * N_i * H
                R = R + N_i * S * (N_i - 1) * H
                D = D + 1500 * N_i * H
        R = R + sum over shared projects of w_p * (V - 1)
        # Firestore cost of day t (free limits reset each day)
        C_day = (0.09*pos(W-20000) + 0.03*pos(R-50000)
                 + 0.01*pos(D-20000)) / 10^5
        total_Firestore = total_Firestore + C_day
        # Claude cost of day t
        T_t = sum over students u of min(X_{u,t}, q)
        n_t = sample from Poisson(mu)
        k_t = ceil(n_t / 40)
        if n_t = 0: C_check = 0
        else:       C_check = 0.116 + 0.106 * (k_t - 1)
        total_Claude = total_Claude + T_t * C_translate + C_check
        total_calls = total_calls + T_t + k_t
    c = total_calls + 1000 * P
    C_otherFirebase = C_storage + C_calls(c) + C_outbound
                      + C_hosting + C_hostingStorage
    C_fixed = scheduler + secrets + backups + PITR + registry
              + builds + Resend + domain
    C_total = total_Claude + total_Firestore + C_otherFirebase
              + C_fixed + C_AppCheck + C_getWebpage + C_dev + C_oneOff
    return C_total and each of its parts
\end{verbatim}

\paragraph{Pseudocode for the full simulation.}
\begin{verbatim}
function MonteCarlo(inputs, n_sim):
    for r = 1 to n_sim:
        C[r] = SimulateMonth(inputs)
    mean = average of C[1..n_sim]
    s    = standard deviation of C[1..n_sim]
    CI   = mean +/- 1.96 * s / sqrt(n_sim)
    return mean, s, CI
\end{verbatim}
The interval $\bar{C} \pm 1.96\, s/\sqrt{n_{\text{sim}}}$ is a 95\% confidence interval for the true expected monthly cost. We choose $n_{\text{sim}}$ large enough (for example 1000) that this interval is narrow, and we report it with every average.

\subsection{Step 3: Analysis}

\paragraph{Cost drivers.}
For each simulated month, we save every part of the cost separately ($C_{\text{Firestore}}$, $C_{\text{Claude}}$, $C_{\text{fixed}}$, and so on). Averaging each part over all months tells us what share of the bill each part makes up.

\paragraph{Sensitivity analysis.}
Many of our inputs are estimates, so we check how much they matter. We change one input at a time (for example, up and down by 20\%), keep the others fixed, rerun the simulation, and record how much the average monthly cost changes. The inputs that move the cost the most are the ones whose real values we most need to know.

\paragraph{Choosing the daily cap.}
The cost never goes down when $q$ goes up, because a higher cap can only allow more translation requests. So we can find the largest cap $q \in [q_{\min}, q_{\max}]$ with average cost at most $B$ by bisection (trying the middle value and cutting the interval in half each time). To make the comparison fair, we reuse the same random numbers for every value of $q$.
\begin{verbatim}
function FindCap(q_min, q_max, B):
    if MonteCarlo(q_max).mean <= B: return q_max
    if MonteCarlo(q_min).mean >  B: return "no cap works"
    while q_max - q_min > 1:
        q = floor((q_min + q_max) / 2)
        if MonteCarlo(q).mean <= B: q_min = q
        else:                       q_max = q
    return q_min
\end{verbatim}

\subsection{Checking the code}

Before trusting the simulation, we test it on cases where we know the answer. If we set all randomness to a fixed value (for example $p_{\text{meet}}=1$ and one class of 30 students with $H=1$), the simulation must give the same Firestore cost as our hand calculation, \$0.115 per day. We also check that the baseline and the simulation agree when the random numbers have no spread, and that the simulated average goes up when we increase $N_{\text{tot}}$.
\end{document}
